% Neuausrichtung der Masterarbeit ab 12.09.2026
% Vorheriger Stand: Archive/Thesis_PreRealignment_2026-09-12/Kapitel/

\section{Theoretischer Hintergrund und Stand der Technik}
\label{sec:grundlagen}

\subsection{MBSE, SysML v2 und maschinenverarbeitbare Modelle}
\label{sec:mbse_sysml}

Systems Engineering adressiert die disziplinübergreifende Entwicklung komplexer
technischer Systeme und verbindet dabei unter anderem Anforderungen,
Architekturentwicklung, Integration sowie Verifikation und Validierung. Model-Based
Systems Engineering (MBSE) erweitert diesen Ansatz durch die systematische Nutzung
von Modellen als zentrale Träger von Engineering-Information. Ziel ist dabei nicht
lediglich, dokumentbasierte Informationen grafisch abzubilden. Vielmehr sollen
fachliche Zusammenhänge und Beziehungen zwischen Engineering-Artefakten explizit
repräsentiert und dadurch konsistenter, nachvollziehbarer und maschinell
verarbeitbar gemacht werden \cite{Hendriks2022}.

Diese Verschiebung von überwiegend dokumentenorientierten zu stärker
modellbasierten Informationsstrukturen bildet eine wesentliche Voraussetzung für
die in dieser Arbeit untersuchte Transformation. Werden Anforderungen, Funktionen,
Strukturen und ihre Beziehungen explizit modelliert, entsteht eine
Informationsrepräsentation, die nicht ausschließlich durch ihre Darstellung,
sondern auch durch ihre definierte Struktur und Semantik verarbeitet werden kann.
MBSE ist dabei nicht als vollständige Ablösung sämtlicher bestehender
Informationsquellen zu verstehen. Gerade in Brownfield-Situationen können
unterschiedliche Dokumente, Datenquellen und Modelle parallel bestehen. Die
Überführung dieser bestehenden Information in einen modellbasierten Kontext wird
in Abschnitt~\ref{sec:brownfield} näher betrachtet.

Als formale Modellierungssprache wird in dieser Arbeit SysML~v2 verwendet.
SysML~v2 wurde für die Beschreibung komplexer Systeme entwickelt und ermöglicht
unter anderem die Repräsentation von Anforderungen, Verhalten, Strukturen,
Use Cases und den Beziehungen zwischen entsprechenden Modellelementen
\cite{SysMLv2}. Die Sprache baut auf KerML als semantischer Grundlage auf und
definiert sowohl grafische als auch textuelle Repräsentationsformen. Damit kann
dieselbe modellierte Engineering-Information in unterschiedlichen Sichten
verwendet werden, ohne dass die grafische Darstellung selbst zum alleinigen
Informationsträger wird.

Für die vorliegende Arbeit ist insbesondere die textuelle Repräsentation von
SysML~v2 relevant. Formal strukturierte Modellinformation kann dadurch als
maschinenverarbeitbares Artefakt erzeugt, gespeichert, verglichen und technisch
geprüft werden. Aktuelle Arbeiten zu SysML~v2 greifen diese Eigenschaften
insbesondere im Zusammenhang mit ausführbaren, überprüfbaren und navigierbaren
Systemmodellen auf \cite{Teodorov2025}. Die textuelle Repräsentation ersetzt
grafische Sichten dabei nicht. Beide dienen unterschiedlichen Formen der
Interaktion mit demselben Modell und können sich insbesondere bei menschlicher
Analyse und automatisierter Verarbeitung ergänzen.

Die Behandlung von Engineering-Artefakten als formal beschriebene und
maschinenverarbeitbare Repräsentationen weist zugleich Bezüge zu
\emph{Everything as Code} auf. Dieses Paradigma überträgt Eigenschaften
codezentrierter Entwicklungspraktiken, beispielsweise textuelle
Repräsentation, Versionsverwaltung und automatisierte Verarbeitung, auf weitere
Entwicklungsartefakte \cite{Stirbu2022}. Für Architekturbeschreibungen wird
dieser Gedanke unter dem Begriff \emph{Architecture as Code} aufgegriffen.
Architektur wird dabei als explizites, versionierbares und automatisiert
verarbeitbares Artefakt behandelt \cite{Bucaioni2025}. Der Begriff wird in
dieser Arbeit nicht mit einer bestimmten Softwarearchitekturmethodik
gleichgesetzt. Relevant ist vielmehr das zugrunde liegende Prinzip, formale
Engineering-Information so zu repräsentieren, dass sie in wiederholbare und
technisch überprüfbare Verarbeitungsabläufe eingebunden werden kann.

SysML~v2 besitzt damit in der vorliegenden Arbeit eine doppelte Rolle. Zum
einen wird die Sprache zur modellbasierten Beschreibung des entwickelten
Architekturkonzepts verwendet. Zum anderen bildet sie die Zielrepräsentation
für Engineering-Information, die aus bestehenden Quellen in einen formalen
Modellkontext überführt werden soll. Diese doppelte Verwendung verbindet die
modellbasierte Architekturentwicklung mit der später untersuchten
KI-gestützten Informationsverarbeitung. Bevor diese näher betrachtet wird,
ist zunächst zu untersuchen, welche besonderen Herausforderungen sich aus der
Transformation bereits bestehender Engineering-Information ergeben.

\subsection{Brownfield Engineering und heterogene Bestandsinformation}
\label{sec:brownfield}

Die Einführung modellbasierter Arbeitsweisen erfolgt in bestehenden
Entwicklungsumgebungen häufig unter Brownfield-Bedingungen. Im Gegensatz zu
einer Greenfield-Situation existieren bereits technische Systeme,
Engineering-Artefakte, historische Entscheidungen und gewachsene
Informationsbestände. Relevante Systeminformation kann dabei in
Spezifikationen, Tabellen, technischen Dokumenten, bestehenden Modellen oder
weiteren Datenquellen verteilt sein. Für eine Transformation zu MBSE besteht
die Aufgabe daher nicht ausschließlich darin, neue Modelle zu erstellen,
sondern vorhandene Engineering-Information zu erschließen und in einen
modellbasierten Zusammenhang zu überführen. Die ineffiziente Übertragung
bestehender Informationen in modellbasierte Repräsentationen wird in diesem
Zusammenhang ausdrücklich als eine Herausforderung der MBSE-Einführung
beschrieben \cite{Maheshwari2018,Henderson2024}.

Die vorhandenen Informationsbestände bilden dabei keine notwendigerweise
einheitliche oder vollständig formalisierte Ausgangsbasis. Quellen können sich
hinsichtlich Format, Strukturierungsgrad, Terminologie, Granularität,
Abstraktionsebene und Aktualität unterscheiden. Zusätzlich können Informationen
über mehrere Artefakte verteilt sein oder nur aus ihrem jeweiligen fachlichen
Kontext verständlich werden. Mehrere Quellen können denselben Systemaspekt aus
unterschiedlichen Perspektiven beschreiben, sich gegenseitig ergänzen oder
voneinander abweichende Aussagen enthalten. Die Heterogenität bestehender
Engineering-Information ist deshalb nicht ausschließlich als technisches
Formatproblem zu betrachten, sondern umfasst auch eine semantische Dimension.

Eine weitere Herausforderung ergibt sich aus der Veränderlichkeit bestehender
Systeme und ihrer Dokumentation. Engineering-Artefakte entwickeln sich mit dem
System weiter und können während seines Lebenszyklus angepasst, erweitert oder
ersetzt werden. Arbeiten zur Co-Evolution technischer Entwicklungsartefakte
zeigen entsprechend die Notwendigkeit, zusammengehörige Engineering-Information
bei Änderungen konsistent weiterzuentwickeln \cite{Legat2014}. Für eine
Brownfield-Transformation bedeutet dies, dass die Beziehung zwischen einer
bestehenden Informationsquelle und daraus abgeleiteter Modellinformation
langfristig relevant bleiben kann. Eine einmalige Übertragung ohne
nachvollziehbare Herkunft erschwert dagegen die spätere Bewertung, welche
Modellanteile durch eine Änderung ihrer Ausgangsinformation betroffen sein
könnten.

Von dieser Problemklasse sind klassische Modelltransformationen zu
unterscheiden. Modellgetriebene Transformationsverfahren überführen
Information zwischen bereits formal beschriebenen Quell- und Zielmodellen
beziehungsweise zwischen definierten Metamodellen. Beispiele hierfür sind
QVT-basierte Transformationen, XMI-basierte Austausch- und
Transformationsverfahren sowie Transformationen zwischen unterschiedlichen
SysML-Repräsentationen \cite{Kapos.2014,Li.2022,Zhou.2024,Meyers.2013}.
Solche Verfahren können auf expliziten Strukturen und definierten semantischen
Beziehungen der Ausgangsmodelle aufbauen.

Bei natürlichsprachlichen oder nur teilweise strukturierten
Bestandsinformationen ist eine solche formale Ausgangsbasis dagegen nicht
vollständig gegeben. Bevor ein Inhalt in eine Zielstruktur überführt werden
kann, muss zunächst bestimmt werden, welche Engineering-Bedeutung er besitzt,
auf welchen Systemgegenstand er sich bezieht und welche Beziehung zu bereits
vorhandener Information besteht. Die Überführung heterogener
Legacy-Information in ein Systemmodell umfasst daher zusätzlich zur formalen
Transformation eine vorgelagerte Interpretations- und Einordnungsaufgabe.

Aus dieser Unterscheidung ergeben sich grundlegende Anforderungen an eine
automatisiert unterstützte Verarbeitung. Eingangsinformation muss so
bereitgestellt werden, dass ihre Identität, Herkunft und ihr relevanter Kontext
während der weiteren Verarbeitung erhalten bleiben können. Gleichzeitig darf
eine fehlende oder nicht eindeutig ableitbare Information nicht allein zur
Vervollständigung einer erwarteten Modellstruktur ergänzt werden. Die Qualität
einer späteren Modellrepräsentation hängt damit nicht nur von deren formaler
Gültigkeit ab, sondern bereits davon, wie die zugrunde liegende
Engineering-Information erschlossen und kontrolliert weiterverarbeitet wird.

Gerade für diese semantische Erschließung bieten Verfahren der
KI-gestützten Sprach- und Informationsverarbeitung neue Möglichkeiten. Ihre
Eignung für Engineering-Anwendungen hängt jedoch davon ab, wie zuverlässig
unterschiedlich strukturierte Quellen interpretiert werden können und wie mit
Unsicherheit, Mehrdeutigkeit und möglichen Fehlinterpretationen umgegangen
wird. Abschnitt~\ref{sec:llm_mbse} betrachtet daher im nächsten Schritt den
Stand KI- und LLM-gestützter Ansätze zur Verarbeitung und Erzeugung von
Engineering-Artefakten.

\subsection{Literaturbasierte Strukturierung der Informationsverarbeitung}
\label{sec:dreischichtenmodell}

Die in Abschnitt~\ref{sec:brownfield} beschriebene Transformation heterogener
Engineering-Bestandsinformation lässt sich nicht auf einen einzelnen
Konvertierungsschritt reduzieren. Bereits die initiale Literaturauswertung
zeigte, dass unterschiedliche Aufgaben der Informationsverarbeitung voneinander
unterschieden werden müssen. Für die weitere Betrachtung wird die
Transformationsaufgabe daher in drei miteinander gekoppelte Ebenen gegliedert:
eine Datenebene, eine Prozessebene und eine Wissensebene. Diese Struktur stellt
kein eigenständiges Referenzmodell aus einer einzelnen Literaturquelle dar,
sondern eine Synthese der für die Problemstellung relevanten Ansätze aus der
initialen Literaturrecherche.

Die \emph{Datenebene} umfasst die vorhandenen Engineering-Informationen und
deren technische sowie strukturelle Eigenschaften. Hierzu gehören unter anderem
unterschiedliche Dateiformate, Strukturierungsgrade, Granularitäten und
Informationsquellen. Im Brownfield-Kontext besteht eine wesentliche
Herausforderung darin, solche heterogenen Informationsbestände zunächst
zugänglich und für eine weitere Verarbeitung eindeutig referenzierbar zu
machen. Arbeiten zur Verarbeitung heterogener Daten- und Wissensbestände zeigen,
dass Unterschiede zwischen Quellen nicht ausschließlich auf deren technisches
Format beschränkt sind, sondern auch unterschiedliche Schemata und
Informationsstrukturen betreffen können
\cite{Maheshwari2018,Tian.2025,Chis.2025}.

Die \emph{Prozessebene} beschreibt die Verarbeitungsschritte, durch die
Information aus ihrem Ausgangszustand in eine formale Zielrepräsentation
überführt wird. Hierzu zählen beispielsweise Transformation, Zuordnung,
Zusammenführung und Validierung von Engineering-Information. Bestehende
modellgetriebene Transformationsansätze zeigen, dass solche Übergänge bei
formal beschriebenen Quell- und Zielmodellen durch explizite Regeln und
Metamodelle gesteuert werden können
\cite{Kapos.2014,Li.2022,Zhou.2024}. Für heterogene Brownfield-Information
reicht eine solche rein strukturelle Transformation jedoch nicht aus, da die
fachliche Bedeutung der Ausgangsinformation zunächst erschlossen werden muss.

Die hierfür erforderliche fachliche Einordnung bildet die
\emph{Wissensebene}. Sie umfasst Begriffe, semantische Beziehungen,
Domänenwissen und Regeln, anhand derer Informationen interpretiert und in einen
gemeinsamen fachlichen Kontext eingeordnet werden können. Ontologien und andere
formalisierte Wissensrepräsentationen werden in diesem Zusammenhang eingesetzt,
um Begriffe und Relationen explizit zu beschreiben und heterogene Informationen
semantisch aufeinander abzubilden
\cite{Abonyi.2024,Meyers.2013,Meng.2025}. Die Wissensebene stellt damit die
Verbindung zwischen den technisch verfügbaren Eingangsdaten und ihrer
fachlichen Bedeutung im späteren Systemmodell her.

Die drei Ebenen sind nicht als voneinander unabhängige Verarbeitungsschichten
zu verstehen. Entscheidungen auf der Wissensebene können beispielsweise
beeinflussen, wie Informationen auf der Prozessebene zusammengeführt werden,
während Eigenschaften der Datenebene begrenzen können, welche
Interpretationen überhaupt belastbar möglich sind. Die Qualität des
resultierenden Modells hängt deshalb nicht ausschließlich von einem einzelnen
Transformationsalgorithmus ab, sondern vom kontrollierten Zusammenspiel von
Informationsgrundlage, Verarbeitung und fachlicher Einordnung.

KI-basierte Verfahren und menschliche Kontrolle werden in dieser Struktur
nicht als zusätzliche vierte Ebene aufgefasst. Large Language Models können
vielmehr innerhalb der beschriebenen Ebenen beziehungsweise an deren
Übergängen Aufgaben übernehmen, für die eine semantische Interpretation
erforderlich ist. Gleichzeitig kann menschliche Entscheidung dort notwendig
werden, wo Informationen mehrdeutig sind oder eine fachliche Autorisierung
erfordern. Bestehende Ansätze zur KI-gestützten Modellgenerierung kombinieren
entsprechend generative Verfahren mit Retrieval, formalen Prüfungen und
menschlicher Einbindung \cite{Cibrian.2025,Topcu2025}.

Das Drei-Ebenen-Modell dient damit als prinzipieller Strukturierungsrahmen für
die betrachtete Informationsverarbeitung. Es beschreibt, \emph{was} innerhalb
einer Brownfield-Transformation berücksichtigt werden muss, ohne bereits eine
konkrete technische Architektur oder Implementierung vorzugeben. Auf dieser
Grundlage wird im folgenden Abschnitt untersucht, welche bestehenden
Mechanismen eingesetzt werden, um KI-gestützte Interpretation innerhalb einer
solchen Verarbeitung kontrollierbar zu gestalten.